\documentclass[11pt,a4paper]{article}

\usepackage[T1]{fontenc}
\usepackage[margin=15mm,top=14mm,bottom=15mm,headheight=13pt,headsep=5mm,footskip=8mm]{geometry}
\usepackage{xcolor}
\usepackage{listings}
\usepackage{enumitem}
\usepackage{tabularx}
\usepackage{fancyhdr}
\usepackage{lastpage}
\usepackage{microtype}
\DisableLigatures{encoding=T1,family=cmtt}

\pagestyle{fancy}
\fancyhf{}
\lhead{\small\textbf{CS 103 -- Quiz 1}}
\rhead{\small IIT Indore -- Autumn 2026}
\cfoot{\small Page \thepage\ of \pageref{LastPage}}
\renewcommand{\headrulewidth}{0.4pt}
\renewcommand{\footrulewidth}{0pt}

\setlength{\parindent}{0pt}
\setlength{\parskip}{3pt}
\setlist[enumerate]{leftmargin=*,itemsep=5pt,topsep=3pt}

\lstdefinestyle{cpp}{
  language=C++,
  basicstyle=\ttfamily\footnotesize,
  frame=single,
  framesep=4pt,
  xleftmargin=2pt,
  xrightmargin=2pt,
  aboveskip=3pt,
  belowskip=3pt,
  showstringspaces=false,
  upquote=true,
  keepspaces=true,
  columns=fullflexible,
  tabsize=4
}

\lstdefinestyle{cppfill}{
  style=cpp,
  basicstyle=\ttfamily\fontsize{20}{24}\selectfont,
  framesep=7pt
}

\newcommand{\answerline}{%
  \par\vspace{4mm}\noindent\rule{\linewidth}{0.35pt}\vspace{1pt}}

\newcommand{\tfcode}[1]{%
  \begingroup
  \setlength{\fboxsep}{1pt}%
  \setlength{\fboxrule}{0.3pt}%
  \fcolorbox{black}{gray!18}{\footnotesize\ttfamily\strut\detokenize{#1}}%
  \endgroup}

\newcommand{\outputbox}[1]{%
  \par\smallskip
  \noindent\textbf{Exact output:}\par
  \vspace{#1}
  \noindent\rule{\linewidth}{0.45pt}%
}

\newcommand{\tfanswer}{%
  \par\noindent
  \textbf{T/F: }\begingroup\setlength{\fboxsep}{1pt}%
  \fbox{\rule{0pt}{4mm}\hspace{15mm}}\endgroup\hspace{4mm}%
  \textbf{Reason: }\rule{0.62\linewidth}{0.35pt}}

\begin{document}

\begin{center}
  {\large\textbf{Indian Institute of Technology Indore}}\\[2pt]
  {\LARGE\textbf{CS 103: Introduction to Programming}}\\[2pt]
  {\Large\textbf{Quiz 1}}
\end{center}

\vspace{2pt}
\begin{tabularx}{\linewidth}{@{}X X r@{}}
  \textbf{Date:} 25 August 2026 &
  \textbf{Time:} 10:35--11:20 (45 minutes) &
  \textbf{Maximum marks:} 42
\end{tabularx}

\vspace{3pt}
\begin{tabularx}{\linewidth}{@{}X X@{}}
  \textbf{Name:} \rule{0.72\linewidth}{0.4pt} &
  \textbf{Roll no.:} \rule{0.62\linewidth}{0.4pt}
\end{tabularx}

\vspace{4pt}
\noindent\fbox{%
  \parbox{\dimexpr\linewidth-2\fboxsep-2\fboxrule\relax}{%
    \textbf{Instructions}
    \begin{enumerate}[label=\arabic*.,leftmargin=5mm,itemsep=1pt,topsep=2pt]
      \item Answer every question in the space provided. Assume C++17 and that all required headers are present.
      \item For True/False, give the verdict and explain why C++ produces that result. Rephrasing the statement is not a reason.
      \item For output questions, reproduce every character, space, and line break exactly. Use an underscore to make a handwritten space visible if needed.
      \item For code completion, write only the missing fragment. It must satisfy the stated goal and terminate.
    \end{enumerate}
  }}

\textbf{Question 1. True or False, with a one-sentence reason.}\hfill
\textbf{12 $\times$ 1 = 12}

Each part needs a correct verdict and an explanation of why C++ produces that result. Merely restating the claim earns no mark.

\begin{enumerate}[label=\textbf{(\alph*)},itemsep=1pt]
  \item A file containing only \tfcode{std::cout << "Hi";} is a complete runnable C++ program.
  \tfanswer

  \item After \tfcode{int x = 9 / 4;}, the variable \tfcode{x} holds \tfcode{2}.
  \tfanswer

  \item After \tfcode{const int limit = 10;}, the statement \tfcode{limit = 12;} is valid because \tfcode{limit} is still an \tfcode{int}.
  \tfanswer

  \item This use of \tfcode{steps} is valid: \tfcode{{ int steps = 3; } std::cout << steps;}.
  \tfanswer

  \item With \tfcode{int x = 8;}, an \tfcode{if (x > 0)}--\tfcode{else if (x > 5)} chain can execute both branches.
  \tfanswer

  \item The expression \tfcode{switch (2.5)} can directly choose a case such as \tfcode{case 2:}.
  \tfanswer

  \item When \tfcode{code} is 1, a \tfcode{case 1} without \tfcode{break} followed by \tfcode{case 2} executes only \tfcode{case 1}.
  \tfanswer

  \item With \tfcode{x = 0}, the condition \tfcode{x != 0 && 12 / x > 2} attempts division by zero.
  \tfanswer

  \item With \tfcode{n = 0}, \tfcode{while (n > 0) { std::cout << n; }} prints \tfcode{0} once.
  \tfanswer

  \item With \tfcode{n = 0}, \tfcode{do { std::cout << n; } while (n > 0);} prints \tfcode{0} once.
  \tfanswer

  \item After reading \tfcode{x}, \tfcode{while (x != -1) { total += x; std::cin >> x; }} adds the sentinel \tfcode{-1} to \tfcode{total}.
  \tfanswer

  \item The loop \tfcode{for (int i = 5; i < 5; i += 1) std::cout << i;} prints \tfcode{5} once.
  \tfanswer
\end{enumerate}
\newpage
\textbf{Question 2. Predict the exact output.}\hfill\textbf{16 marks}

Parts (a)--(e) carry 2 marks each; parts (f)--(g) carry 3 marks each.

\noindent\begin{minipage}[t]{0.485\linewidth}
\textbf{(a)}
\begin{lstlisting}[style=cpp]
int distance = 7, time = 2;
double speed = distance / time;
std::cout << speed << '\n';
\end{lstlisting}
\outputbox{7mm}
\end{minipage}\hfill
\begin{minipage}[t]{0.485\linewidth}
\textbf{(b)}
\begin{lstlisting}[style=cpp]
int score = 7;
if (score > 7) std::cout << "A ";
else if (score >= 5) std::cout << "B ";
else std::cout << "C ";
std::cout << score % 4 << '\n';
\end{lstlisting}
\outputbox{7mm}
\end{minipage}

\vspace{4pt}
\noindent\begin{minipage}[t]{0.485\linewidth}
\textbf{(c)}
\begin{lstlisting}[style=cpp]
int n = 3;
while (n > 0) {
    std::cout << n << ' ';
    n -= 1;
}
\end{lstlisting}
\outputbox{7mm}
\end{minipage}\hfill
\begin{minipage}[t]{0.485\linewidth}
\textbf{(d)}
\begin{lstlisting}[style=cpp]
int total = 0;
for (int i = 1; i <= 6; i += 1) {
    if (i % 2 == 0) continue;
    total += i;
}
std::cout << total << '\n';
\end{lstlisting}
\outputbox{7mm}
\end{minipage}

\vspace{4pt}
\noindent\begin{minipage}[t]{0.485\linewidth}
\textbf{(e)}
\begin{lstlisting}[style=cpp]
int code = 1;
switch (code) {
 case 0: std::cout << "zero "; break;
 case 1: std::cout << "one ";
 case 2: std::cout << "two "; break;
 default: std::cout << "other ";
}
\end{lstlisting}
\outputbox{7mm}
\end{minipage}\hfill
\begin{minipage}[t]{0.485\linewidth}
\textbf{(f)}
\begin{lstlisting}[style=cpp]
int total = 0;
for (int r = 1; r <= 3; r += 1) {
    for (int c = 1; c <= r; c += 1) {
        if (c == 2) continue;
        total += r + c;
    }
    std::cout << total << ' ';
}
std::cout << '\n';
\end{lstlisting}
\outputbox{7mm}
\end{minipage}

\vspace{4pt}
\noindent\textbf{(g)}
\begin{lstlisting}[style=cpp]
int x = 0;
if (x != 0 && 12 / x > 2)
    std::cout << "A ";
else
    std::cout << "B ";

switch (x + 1) {
    case 0: std::cout << "zero ";
    case 1: std::cout << "one ";
    case 2: std::cout << "two "; break;
    default: std::cout << "other ";
}
std::cout << '\n';
\end{lstlisting}
\outputbox{7mm}

\newpage
\textbf{Question 3. Complete the missing code.}\hfill\textbf{14 marks}

Write each answer directly in its code blank. Parts (a)--(d) carry 2 marks each; parts (e)--(f) carry 3 marks each.

\textbf{(a)} Stop at the sentinel \lstinline|0| and do not add it.
\begin{lstlisting}[style=cppfill]
int value = 0, total = 0;
std::cin >> value;
while (__________________) {
    total += value;
    std::cin >> value;
}
std::cout << total;
\end{lstlisting}

\vspace{5pt}
\textbf{(b)} Keep asking until \lstinline|choice| is from 1 through 4.
\begin{lstlisting}[style=cppfill]
int choice = 0;
do {
    std::cin >> choice;
} while (____________________);
\end{lstlisting}

\vspace{5pt}
\textbf{(c)} Print \lstinline|2 4 6 8 10|, with a space after each.
\begin{lstlisting}[style=cppfill]
for (int k = 2; k <= 10; ________) {
    std::cout << k << ' ';
}
\end{lstlisting}

\vspace{5pt}
\textbf{(d)} Run the division only for a nonzero denominator.
\begin{lstlisting}[style=cppfill]
if (__________________ &&
    numerator / denominator > 2) {
    std::cout << "large quotient";
}
\end{lstlisting}

\newpage
\textbf{Question 3. Complete the missing code (continued).}\hfill\textbf{14 marks}

\begin{enumerate}[label=\textbf{(\alph*)},start=5,itemsep=8pt]
\item The input is a sequence of integers ending with \lstinline|-1|. Complete the three blanks so the program prints the integer average of the \emph{positive even} values before the sentinel, or \lstinline|none| if there are no such values. \hfill\textbf{3 marks}

\begin{lstlisting}[style=cppfill]
int x = 0;
int total = 0;
int count = 0;
std::cin >> x;
while (__________________) {
    if (__________________) {
        total += x;
        __________________;
    }
    std::cin >> x;
}
if (count == 0)
    std::cout << "none\n";
else
    std::cout << total / count << '\n';
\end{lstlisting}

\item Complete all three blanks so the program prints \lstinline|1| if the input integer $n$ has a divisor from 2 through $n-1$, and \lstinline|0| otherwise. \hfill\textbf{3 marks}
\begin{lstlisting}[style=cppfill]
int n = 0;
std::cin >> n;
bool hasDivisor = false;
for (int d = 2; __________; d += 1) {
    if (________________) {
        ________________;
        break;
    }
}
std::cout << hasDivisor;
\end{lstlisting}
\end{enumerate}

\end{document}
